\section*{Erd\H{o}s Problem 905: an edge in many triangles}

\paragraph{Statement and result.}
Every \(n\)-vertex graph with \(m>n^2/4\) edges has an
edge belonging to more than \(n/6\) triangles.
We prove this by elementary counting, in the form of
the booksize inequality used by Bollob\'as and Nikiforov.

\paragraph{Notation and two local bounds.}
For an edge \(uv\), let
\(t_{uv}=|N(u)\cap N(v)|\), and put
\[
 b=\max_{uv\in E}t_{uv},\qquad T=\text{number of triangles}.
\]
Then \(\sum_{uv\in E}t_{uv}=3T\).
Since \(|N(u)\cup N(v)|\leq n\), every edge satisfies
\[
 d(u)+d(v)-n\leq t_{uv}.
\tag{1}
\]
For a triangle \(uvw\), inclusion--exclusion gives
\[
 d(u)+d(v)+d(w)
 \leq n+t_{uv}+t_{uw}+t_{vw}.
\]
Summing over triangles yields
\[
 \sum_{uv\in E}t_{uv}(d(u)+d(v))
 \leq2nT+2\sum_{uv\in E}t_{uv}^2.
\tag{2}
\]
Indeed each triangle's degree sum occurs twice on the
left before summing, whereas an edge contributes its
codegree once for each of its \(t_{uv}\) triangles on
the right.

\paragraph{A global booksize inequality.}
Set
\[
 S=\sum_{uv\in E}(d(u)+d(v)-n)
   =\sum_vd(v)^2-nm.
\]
Multiply (1) by the nonnegative number \(b-t_{uv}\),
and sum. Using (2) gives
\[
 \begin{split}
 bS
 &\leq\sum_{uv\in E}
       \{t_{uv}(d(u)+d(v)-n)+(b-t_{uv})t_{uv}\}\\
 &=\sum_{uv\in E}t_{uv}(d(u)+d(v))
       +3(b-n)T-\sum_{uv\in E}t_{uv}^2\\
 &\leq(3b-n)T+\sum_{uv\in E}t_{uv}^2\\
 &\leq(6b-n)T.
 \end{split}
\tag{3}
\]
All terms and directions are valid even when an individual
edge has \(d(u)+d(v)-n<0\).

\paragraph{The density assumption forces a large book.}
Cauchy--Schwarz and \(2m=\sum_vd(v)\) give
\[
 S\geq4m^2/n-nm=m(4m/n-n)>0.
\]
If \(b=0\), (1) would give \(S\leq0\), a contradiction.
Thus \(b>0\), and consequently \(T>0\).
Equation (3) now has positive left-hand side, so
\(6b-n>0\). An edge attaining \(b\) is the required edge.

\paragraph{Attribution and scope.}
Edwards and Khadzhiivanov--Nikiforov independently
established the result. A primary exposition of its
booksize formulation is B. Bollob\'as and V. Nikiforov,
\emph{Books in graphs}, \url{https://arxiv.org/abs/math/0405080}.
See also \url{https://www.erdosproblems.com/905}, checked
24 September 2026. The displayed counting proof is complete
and uses no external extremal-graph theorem. It is a
reconstruction of a known result; no local Lean verification
is claimed.

\paragraph{Completion estimate.}
\textbf{COMPLETION: 100\%}
